\subsection{Composition constitutive et évaluation de la section gravitationnelle}
\label{g26:subsec:evaluacion-viii-es}

L'évaluation réutilise les coordonnées régionales générées et le registre
\(K\) sélectionné dans le noyau précédent. La coordonnée \(\alpha\)
détermine les sections d'action de
\eqref{v:ant:eq:el-secciones-accion} ; dans la section de retour, nous écrivons
\(\eta=\eta_{\rm ret}\) et \(S_*=10^{-34}\,\mathrm{J\,s}\), de sorte
que \(\hbar_{\rm ret}=S_*\eta\). L'angle et le contre-angle déjà
construits fournissent, dans la représentation en degrés,
\[
 A_{\deg}=1000\alpha,\qquad C^*_{\deg}=2(\eta+\alpha),\qquad
 x=\frac{\pi_{\rm HMT}A_{\deg}}{180},\quad
 y=\frac{\pi_{\rm HMT}C^*_{\deg}}{180}.
\]
Ainsi, \(\eta=(C^*_{\deg}-2\alpha)/2\). La différence est prise entre
des coordonnées de la même représentation angulaire.

Dans le domaine généré \(0<y<x\), les canaux étiquetés sont
\[
 q_\pm=e^{-(x\pm y)},\qquad
 r_\pm=\frac{1-q_\pm^{90}}{1-q_\pm^{120}},\qquad
 0<q_\pm<1.
\]
Le quotient procède des deux secteurs d'extension : si
\(S_n(q)=\sum_{j=0}^{n-1}q^j\), l'identité
\((1-q)S_n(q)=1-q^n\) démontre
\(r_\pm=S_{90}(q_\pm)/S_{120}(q_\pm)\).
Les carrés des canaux déterminent les lectures constitutives
\[
 \widehat\varepsilon=r_+^2,\qquad
 \widehat\mu=r_-^2,\qquad
 \widehat c=\frac1{r_+r_-},\qquad
 \widehat Z=\frac{r_-}{r_+}.
\]
Pour des bases positives de vitesse \(c_*\), d'action \(u_S\) et de charge
\(u_Q\), les sections dimensionnelles s'expriment sous la forme
\[
 \mu=\widehat\mu\frac{u_S}{c_*u_Q^2},\qquad
 \varepsilon=\widehat\varepsilon\frac{u_Q^2}{u_Sc_*},\qquad
 c=c_*\widehat c=\frac1{\sqrt{\mu\varepsilon}}.
\]
En effet, le produit des deux premières équations élimine les bases
d'action et de charge et laisse
\(\mu\varepsilon=(r_+r_-)^2/c_*^2\) ; sa racine positive prouve la
dernière égalité. Le facteur \(\widehat c\) transforme la base \(c_*\)
une seule fois. L'évaluation SI qui suit représente le même \(c\)
par \(299792458\,\mathrm{m\,s^{-1}}\) et restitue
\(c_*=c/\widehat c\). Cette opération vérifie l'équivalence des
représentations dimensionnelle et constitutive.

Avec \(r_N=5759/23040\) et la référence déclarée \(L_*=1\,\mathrm m\),
la substitution des équations précédentes dans le coefficient radial donne
\begin{align}
 G&=\frac{c^3L_*^2}{S_*}\frac{r_N^2\alpha^{32}}{\eta}
   =\frac{2c^3L_*^2r_N^2\alpha^{32}}
           {S_*(C^*_{\deg}-2\alpha)},\label{g26:eq:g-angular-viii-es}\\
  &=\frac{L_*^2r_N^2\alpha^{32}}
           {S_*\eta(\mu\varepsilon)^{3/2}}.
 \label{g26:eq:g-constitutivo-viii-es}
\end{align}
Ce sont des expressions du même coefficient dans la réalisation radiale. La
référence longitudinale \(L_*\), la longueur construite \(L_\alpha\),
l'arête chronologique \(ct_0\) et le rayon d'horizon conservent leurs
types différenciés.

L'évaluation avec les développements régionaux archivés fournit
\begin{align*}
 L_\alpha&=1.6162549826251263301682625013\ldots\,10^{-35}\,\mathrm m,\\
 \hbar_{\rm ret}&=1.0545718169150390611336905847\ldots\,10^{-34}\,
                    \mathrm{J\,s},\\
 t_0&=3.1364999625365490734\ldots\,10^{-45}\,\mathrm s,\\
 G&=6.6742996594140227063677761891\ldots\,10^{-11}\,
                    \mathrm{m^3\,kg^{-1}\,s^{-2}}.
\end{align*}
Les calculs à 100 et 140 chiffres de travail conservent 71 chiffres significatifs de l'évaluation. Cette stabilité numérique se distingue de l'incertitude expérimentale. Par rapport à la référence CODATA 2022 \(G_{\rm ref}=6.67430(15)\,10^{-11}\,\mathrm{m^3\,kg^{-1}\,s^{-2}}\), la différence est \(G-G_{\rm ref}=-3.4058597729\ldots\,10^{-18}\,
\mathrm{m^3\,kg^{-1}\,s^{-2}}\), soit environ \(-0.00227057\) fois son incertitude-type. La référence est comparée après le calcul ; les règles R4, R5 et l'identification R8 conservent la portée démontrée dans la section précédente.

Les sources de reproduction comprennent les développements régionaux, le
registre de \(K\), les équations d'action et de contre-angle, les
vérificateurs exacts et l'évaluation décimale. Ils s'exécutent à partir du
dossier local \path{reproduccion_gravitatoria_20260926} au moyen de
\path{verificar_reproduccion.py}. Chaque preuve déclare son domaine ; la
réutilisation de coordonnées archivées est distinguée d'une nouvelle
exécution intégrale du générateur.
